\Abstract{The centenary Hilbert-Pólya conjecture proposes that the non-trivial zeros of the Riemann zeta function correspond to eigenvalues of a Hermitian operator. Despite the statistical correlations established by Dyson-Montgomery and computationally verified by Odlyzko across more than $10^{13}$ zeros, the explicit construction of the required operator has resisted all attempts due to fundamental constructional circularity.

We approach the problem departing from the perspective of Manin's Numbers as Functions and $\mathbb{F}_1$-geometric program [23,24], where $\text{Spec}(\mathbb{Z})$ is treated as a geometric object over the field with one element, and the program's established connection to string theory via Connes, Douglas and Schwarz (1998), who demonstrated that toroidal compactification in M-theory produces the noncommutative tori central to this program [27]. While this framework constructs the geometric arena an arithmetic proof requires, it lacks a dynamical component. We address this gap by constructing a Hermitian operator $T^*$ on a toroidal manifold whose spectral invariants ($d = 3$, $\mu = 1/2$, $\sigma = 3/2$) emerge from the generator matrix $\hat{\Omega} = \frac{1}{2} \cdot \text{diag}(1,\omega,\omega^2)$ and whose structural parameters derive from geometric-arithmetic architecture --- hypercube stratification, Mersenne boundaries, and Golden Prime classes (Grisales Herrera, 2025) --- without utilizing known values of the zeros.

Computational verification yields mean error 0.59\% for $n = 50$--$100$ and 0.25\% for $n = 1,000$--$10,000$, with the ratio $T^*(n)/t_n$ converging toward unity, validated to $n = 131,072$.

This constitutes, to our knowledge, the first non-circular construction of a Hilbert-Pólya operator with verified asymptotic convergence $T^*(n)/t_n \to 1$, suggesting that the Riemann Hypothesis may admit reinterpretation as a geometric stability condition within the $\mathbb{F}_1$-arithmetic framework and that the emergent identity $\zeta(2)/(\pi/3)^2 = \sigma$, connecting the Euler product with the $S_3$ angular structure, suggests the mechanism has scope beyond the Riemann spectrum.}

\Keywords{Riemann Hypothesis; Hilbert-Pólya conjecture; Hermitian operators; Non-circular construction; $S_3$ symmetry; Spectral invariants; Golden ratio; Fibonacci sequence}

\Classification{11M26; 11G05; 81Q10; 81T30}
